\begin{recipe}
[% 
    preparationtime = {\unit[10]{min}},
    bakingtime={\unit[35]{min}},
    bakingtemperature={\protect\bakingtemperature{
        topbottomheat=\unit[180]{\textcelcius}}},
    portion = {\portion{16}},
    source = {Malena Barillas Dahm},
    calory = {200 kcal | 3g Protein | 30g KH | 8g Fett},
    rating = {3}
]
{Fudgy Vegane Brownies}
    
    % \graph
    % {% pictures
    %     %small=pic/brownie_small,
    %     %big=pic/brownie_big
    % }
    
    \introduction{%
        Schnell und einfach gemacht, wenn er noch leicht fudgy ist und nicht durchgebacken merkt man nicht die fehlenden Eier.
    }
    
    \ingredients{%
        1\nicefrac{1}{2} Cup & Brauner Zucker\\
        \nicefrac{1}{2} Cup & Wasser\\
        5 EL & Apfelmus\\
        \nicefrac{1}{2} Cup & Neutrales Öl\\
        \nicefrac{3}{4} Cup & Kakaopulver\\
        1\nicefrac{1}{2} Cup & Mehl\\
        \nicefrac{1}{2} TL & Salz\\
        Optional & Schokodrops, Himbeeren
    }
    
    \preparation{%
        \begin{enumerate}
            \item Backofen auf \unit[180]{\textcelcius} vorheizen.
            
            \item Zucker, Wasser, Apfelmus und Öl in einer großen Schüssel gut verrühren.
            
            \item Kakaopulver, Mehl und Salz hinzufügen und alles zu einem glatten Teig vermengen. Optional Schokodrops oder Himbeeren unterheben.
            
            \item Eine Form (ca. 20 x 30 cm) einfetten oder mit Backpapier auslegen und den Teig gleichmäßig verteilen.
            
            \item Etwa 35 Minuten backen. Mit einer Gabel oder einem Holzstäbchen testen, ob der Brownie fertig ist - innen darf er noch leicht fudgy sein.
            
            \item Vor dem Anschneiden vollständig abkühlen lassen.
        \end{enumerate}
    }
    
    \suggestion{%
        Mit einer Prise Meersalz oder frischen Himbeeren servieren.
    }
    
    \hint{%
        Je kürzer die Backzeit, desto fudgier wird der Brownie. Lieber etwas früher aus dem Ofen nehmen - er zieht beim Abkühlen noch nach.
    }
    
\end{recipe}